\tituloI{METODOLOGÍA}

\tituloII{Método, tipo y alcance de la investigación}

\tituloIII{Método de investigación}
\parrafoIII{%
Se adopta el método cuasi-experimental. Se manipulará la lógica de validación de datos y la arquitectura de comunicación fiscal (variable independiente) para observar su efecto sobre la integridad del inventario y la eficiencia en la emisión de comprobantes (variable dependiente).%
}

\tituloIII{Tipo de investigación}
\parrafoIII{%
Investigación tecnológica y aplicada, orientada a desarrollar un prototipo funcional de sistema de gestión local que valide la hipótesis de mejora en la consistencia de datos y autonomía tributaria en pequeños comercios.%
}

\tituloIII{Alcance de la investigación}
\parrafoIII{%
Alcance descriptivo y correlacional, caracterizando las deficiencias del flujo de trabajo actual en el negocio de abarrotes y estableciendo la relación entre la implementación de reglas de negocio estrictas y la reducción de errores operativos.%
}

\tituloIII{Diseño de la investigación}
\parrafoIII{%
Diseño pre-experimental con pre-prueba y post-prueba para un solo grupo. Se evaluará el impacto de la implementación mediante una comparación entre el estado inicial de los registros de stock (pre-prueba) y los resultados obtenidos tras el uso del sistema propuesto (post-prueba) en condiciones reales de operación durante un periodo de tres meses.%
}

\tituloII{Materiales y métodos}

\tituloIII{Materiales}
\parrafoIII{%
Para la ejecución se utilizarán los siguientes recursos:
}

\begin{itemize}
    \item \textbf{Hardware:} Estación de trabajo con arquitectura x86\_64, 8GB RAM, procesador de 2 núcleos y almacenamiento SSD (configuración representativa de un servidor local de bajo costo).
    
    \item \textbf{Software base:} Arch Linux (entorno minimalista para optimización de recursos) y MariaDB 10.x para la persistencia de datos.
    
    \item \textbf{Lenguajes y librerías:} Python 3.12, framework para API local, y librerías de firma digital (como xmlsig o similares) para la generación de archivos UBL.
    
    \item \textbf{Herramientas de medición:} 
    \begin{itemize}
        \item \textit{Unittest / Pytest}: Para medir la tasa de éxito de las validaciones de stock negativo y reglas de integridad.
        \item \textit{Logs de transacciones}: Para cuantificar el tiempo de respuesta entre la venta y la confirmación del CPE (Comprobante de Pago Electrónico).
    \end{itemize}
    
    \item \textbf{Sistema de Gestión Local:} Prototipo con módulos de inventario, ventas y facturación directa.
\end{itemize}

\tituloIII{Población y muestra}
\parrafoIII{%
La población está compuesta por la totalidad de transacciones diarias y artículos del inventario del local de abarrotes. La muestra incluye un conjunto de 500 operaciones de venta y 50 categorías de productos seleccionadas mediante muestreo por conveniencia. Se realizarán 3 mediciones diarias de consistencia de stock para garantizar la validez estadística de los datos recogidos.%
}

\tituloIII{Variables de estudio}
\begin{itemize}
    \item \textbf{Variables Independientes:}
        \begin{itemize}
            \item Sistema de gestión local: Implementación de algoritmos de validación de existencias y firma digital directa.
        \end{itemize}
    \item \textbf{Variables Dependientes:}
        \begin{itemize}
            \item Integridad del inventario: Cantidad de registros con stock negativo o discrepante.
            \item Autonomía operativa: Reducción del tiempo de dependencia de software externo de terceros.
        \end{itemize}
    \item \textbf{Variables Controladas / Intervinientes:}
        \begin{itemize}
            \item Normativa SUNAT: Se mantiene estática bajo los estándares UBL vigentes.
            \item Estabilidad eléctrica: Uso de sistemas de protección para evitar corrupción de la base de datos local.
        \end{itemize}
\end{itemize}

\tituloIII{Técnicas e instrumentos}
\parrafoIII{%
La técnica empleada es la observación estructurada y la experimentación controlada. El instrumento consiste en una suite automatizada de pruebas de integración y una guía de observación técnica que permitirá contrastar el stock físico frente al stock reportado por el sistema de forma diaria.%
}

\tituloIII{Procedimiento}
\parrafoIII{%
El procedimiento se estructurará en 6 fases:
}

\begin{itemize}
    \item \textbf{Fase 1: Diagnóstico técnico:} Análisis de las causas de la inconsistencia de datos en el sistema actual del negocio.
    \item \textbf{Fase 2: Modelado de datos:} Diseño de la base de datos con restricciones (constraints) de nivel de motor para impedir valores negativos.
    \item \textbf{Fase 3: Desarrollo del core:} Codificación de la lógica de negocio y el módulo de comunicación con el OSE/SUNAT.
    \item \textbf{Fase 4: Integración local:} Instalación del software en el hardware Bare Metal y configuración del entorno Arch Linux.
    \item \textbf{Fase 5: Pruebas de campo:} Ejecución del sistema en paralelo con la operación diaria para verificar la exactitud de los resultados.
    \item \textbf{Fase 6: Evaluación y ajuste:} Refinamiento de las interfaces y optimización de las consultas SQL según los reportes de uso.
\end{itemize}

\tituloIII{Análisis de datos}
\parrafoIII{%
Se aplicará estadística descriptiva para analizar la frecuencia de errores de stock y los tiempos promedio de facturación. Se utilizará el análisis de tendencias para comparar el rendimiento del sistema propuesto frente al método anterior. El procesamiento se realizará con Python y la librería Pandas para la generación de cuadros comparativos y métricas de eficiencia.%
}